% Factsheet Anspruch: Seite 1 Deutsch, Seite 2 Englisch, Abschnitt für Abschnitt
% gleich. Keine Logos, keine Bilder.
% Quellen: docs/proposal-organisations.md (Fakten, was der Partner bekommt, was
% Anspruch nicht ist, Datenschutz), docs/PARTNERS.md, README.md,
% lib/legal/datenschutz.ts.
% Lokale Makros: \fillin (Kasten mit den Platzhaltern).
\documentclass[fontsize=10pt,paper=a4,parskip=half]{scrartcl}
% Shared preamble for the Anspruch dossier and its companion documents.
% Load after \documentclass with \input{preamble} (from docs/dossier/)
% or % Shared preamble for the Anspruch dossier and its companion documents.
% Load after \documentclass with \input{preamble} (from docs/dossier/)
% or % Shared preamble for the Anspruch dossier and its companion documents.
% Load after \documentclass with \input{preamble} (from docs/dossier/)
% or \input{../preamble} (from docs/dossier/antrag/ or docs/dossier/partner/).
% Each document sets its own \documentclass and its own pdftitle via
% \hypersetup{pdftitle={...}}. Only packages that ship with MiKTeX.

\usepackage[T1]{fontenc}
\usepackage[utf8]{inputenc}
\usepackage{lmodern}
\usepackage[english,ngerman]{babel} % last option = main language: ngerman
\usepackage[a4paper,margin=25mm]{geometry}
\usepackage{booktabs}
\usepackage{longtable}
\usepackage{array}
\usepackage{enumitem}
\usepackage{xcolor}
\usepackage{microtype}
\usepackage[autostyle=true,german=quotes]{csquotes}
\usepackage{marginnote}
\usepackage[hidelinks,unicode,pdflang={de-DE}]{hyperref}

% ---------------------------------------------------------------------------
% Colours
\definecolor{blankbg}{RGB}{255,236,153}   % highlighted placeholder
\definecolor{todofg}{RGB}{170,40,40}      % "to confirm" margin note

% \blank{Text}: a highlighted placeholder "[Text]" for values the sender
% fills in before sending. Usable in running text and in table cells.
\newcommand{\blank}[1]{\begingroup\setlength{\fboxsep}{1pt}%
  \colorbox{blankbg}{\strut[#1]}\endgroup}

% \todo{Text}: a "zu bestätigen" note in the margin. Works in running text,
% lists and longtable cells (via marginnote, not \marginpar).
% Defined with \providecommand so that it does not clash if a document
% loads todonotes on its own (todonotes then wins).
\providecommand{\todo}[1]{\marginnote{\raggedright\scriptsize
  \textcolor{todofg}{\textbf{Zu bestätigen:} #1}}}
\setlength{\marginparwidth}{20mm}

% \draftmark: header text for every page of a draft.
\newcommand{\draftmark}{Entwurf, 23.10.2026}

% Draft header for scrartcl / scrreprt (KOMA) and article / report.
% Letters (scrlttr2, letter) place \draftmark themselves.
\makeatletter
\@ifclassloaded{scrartcl}{\def\anspruch@koma{}}{}
\@ifclassloaded{scrreprt}{\def\anspruch@koma{}}{}
\@ifclassloaded{article}{\def\anspruch@std{}}{}
\@ifclassloaded{report}{\def\anspruch@std{}}{}
\ifdefined\anspruch@koma
  \usepackage[headsepline]{scrlayer-scrpage}
  \clearpairofpagestyles
  \ohead{\small\draftmark}
  \ihead{\small\headmark}
  \cfoot*{\pagemark}
  \automark[section]{section}
  \pagestyle{scrheadings}
\fi
\ifdefined\anspruch@std
  \usepackage{fancyhdr}
  \pagestyle{fancy}
  \fancyhf{}
  \fancyhead[R]{\small\draftmark}
  \fancyhead[L]{\small\nouppercase{\leftmark}}
  \fancyfoot[C]{\thepage}
  \fancypagestyle{plain}{\fancyhf{}\fancyhead[R]{\small\draftmark}%
    \fancyfoot[C]{\thepage}\renewcommand{\headrulewidth}{0pt}}
\fi
\makeatother

% Lists: compact, as in a German Antrag.
\setlist{itemsep=0.2em,topsep=0.4em}

% Tables: a bit more air, ragged-right paragraph columns.
\renewcommand{\arraystretch}{1.15}
\newcolumntype{L}[1]{>{\raggedright\arraybackslash}p{#1}}
 (from docs/dossier/antrag/ or docs/dossier/partner/).
% Each document sets its own \documentclass and its own pdftitle via
% \hypersetup{pdftitle={...}}. Only packages that ship with MiKTeX.

\usepackage[T1]{fontenc}
\usepackage[utf8]{inputenc}
\usepackage{lmodern}
\usepackage[english,ngerman]{babel} % last option = main language: ngerman
\usepackage[a4paper,margin=25mm]{geometry}
\usepackage{booktabs}
\usepackage{longtable}
\usepackage{array}
\usepackage{enumitem}
\usepackage{xcolor}
\usepackage{microtype}
\usepackage[autostyle=true,german=quotes]{csquotes}
\usepackage{marginnote}
\usepackage[hidelinks,unicode,pdflang={de-DE}]{hyperref}

% ---------------------------------------------------------------------------
% Colours
\definecolor{blankbg}{RGB}{255,236,153}   % highlighted placeholder
\definecolor{todofg}{RGB}{170,40,40}      % "to confirm" margin note

% \blank{Text}: a highlighted placeholder "[Text]" for values the sender
% fills in before sending. Usable in running text and in table cells.
\newcommand{\blank}[1]{\begingroup\setlength{\fboxsep}{1pt}%
  \colorbox{blankbg}{\strut[#1]}\endgroup}

% \todo{Text}: a "zu bestätigen" note in the margin. Works in running text,
% lists and longtable cells (via marginnote, not \marginpar).
% Defined with \providecommand so that it does not clash if a document
% loads todonotes on its own (todonotes then wins).
\providecommand{\todo}[1]{\marginnote{\raggedright\scriptsize
  \textcolor{todofg}{\textbf{Zu bestätigen:} #1}}}
\setlength{\marginparwidth}{20mm}

% \draftmark: header text for every page of a draft.
\newcommand{\draftmark}{Entwurf, 23.10.2026}

% Draft header for scrartcl / scrreprt (KOMA) and article / report.
% Letters (scrlttr2, letter) place \draftmark themselves.
\makeatletter
\@ifclassloaded{scrartcl}{\def\anspruch@koma{}}{}
\@ifclassloaded{scrreprt}{\def\anspruch@koma{}}{}
\@ifclassloaded{article}{\def\anspruch@std{}}{}
\@ifclassloaded{report}{\def\anspruch@std{}}{}
\ifdefined\anspruch@koma
  \usepackage[headsepline]{scrlayer-scrpage}
  \clearpairofpagestyles
  \ohead{\small\draftmark}
  \ihead{\small\headmark}
  \cfoot*{\pagemark}
  \automark[section]{section}
  \pagestyle{scrheadings}
\fi
\ifdefined\anspruch@std
  \usepackage{fancyhdr}
  \pagestyle{fancy}
  \fancyhf{}
  \fancyhead[R]{\small\draftmark}
  \fancyhead[L]{\small\nouppercase{\leftmark}}
  \fancyfoot[C]{\thepage}
  \fancypagestyle{plain}{\fancyhf{}\fancyhead[R]{\small\draftmark}%
    \fancyfoot[C]{\thepage}\renewcommand{\headrulewidth}{0pt}}
\fi
\makeatother

% Lists: compact, as in a German Antrag.
\setlist{itemsep=0.2em,topsep=0.4em}

% Tables: a bit more air, ragged-right paragraph columns.
\renewcommand{\arraystretch}{1.15}
\newcolumntype{L}[1]{>{\raggedright\arraybackslash}p{#1}}
 (from docs/dossier/antrag/ or docs/dossier/partner/).
% Each document sets its own \documentclass and its own pdftitle via
% \hypersetup{pdftitle={...}}. Only packages that ship with MiKTeX.

\usepackage[T1]{fontenc}
\usepackage[utf8]{inputenc}
\usepackage{lmodern}
\usepackage[english,ngerman]{babel} % last option = main language: ngerman
\usepackage[a4paper,margin=25mm]{geometry}
\usepackage{booktabs}
\usepackage{longtable}
\usepackage{array}
\usepackage{enumitem}
\usepackage{xcolor}
\usepackage{microtype}
\usepackage[autostyle=true,german=quotes]{csquotes}
\usepackage{marginnote}
\usepackage[hidelinks,unicode,pdflang={de-DE}]{hyperref}

% ---------------------------------------------------------------------------
% Colours
\definecolor{blankbg}{RGB}{255,236,153}   % highlighted placeholder
\definecolor{todofg}{RGB}{170,40,40}      % "to confirm" margin note

% \blank{Text}: a highlighted placeholder "[Text]" for values the sender
% fills in before sending. Usable in running text and in table cells.
\newcommand{\blank}[1]{\begingroup\setlength{\fboxsep}{1pt}%
  \colorbox{blankbg}{\strut[#1]}\endgroup}

% \todo{Text}: a "zu bestätigen" note in the margin. Works in running text,
% lists and longtable cells (via marginnote, not \marginpar).
% Defined with \providecommand so that it does not clash if a document
% loads todonotes on its own (todonotes then wins).
\providecommand{\todo}[1]{\marginnote{\raggedright\scriptsize
  \textcolor{todofg}{\textbf{Zu bestätigen:} #1}}}
\setlength{\marginparwidth}{20mm}

% \draftmark: header text for every page of a draft.
\newcommand{\draftmark}{Entwurf, 23.10.2026}

% Draft header for scrartcl / scrreprt (KOMA) and article / report.
% Letters (scrlttr2, letter) place \draftmark themselves.
\makeatletter
\@ifclassloaded{scrartcl}{\def\anspruch@koma{}}{}
\@ifclassloaded{scrreprt}{\def\anspruch@koma{}}{}
\@ifclassloaded{article}{\def\anspruch@std{}}{}
\@ifclassloaded{report}{\def\anspruch@std{}}{}
\ifdefined\anspruch@koma
  \usepackage[headsepline]{scrlayer-scrpage}
  \clearpairofpagestyles
  \ohead{\small\draftmark}
  \ihead{\small\headmark}
  \cfoot*{\pagemark}
  \automark[section]{section}
  \pagestyle{scrheadings}
\fi
\ifdefined\anspruch@std
  \usepackage{fancyhdr}
  \pagestyle{fancy}
  \fancyhf{}
  \fancyhead[R]{\small\draftmark}
  \fancyhead[L]{\small\nouppercase{\leftmark}}
  \fancyfoot[C]{\thepage}
  \fancypagestyle{plain}{\fancyhf{}\fancyhead[R]{\small\draftmark}%
    \fancyfoot[C]{\thepage}\renewcommand{\headrulewidth}{0pt}}
\fi
\makeatother

% Lists: compact, as in a German Antrag.
\setlist{itemsep=0.2em,topsep=0.4em}

% Tables: a bit more air, ragged-right paragraph columns.
\renewcommand{\arraystretch}{1.15}
\newcolumntype{L}[1]{>{\raggedright\arraybackslash}p{#1}}

\hypersetup{pdftitle={Anspruch: Factsheet}}
\RedeclareSectionCommand[beforeskip=1.4ex plus .2ex,afterskip=.3ex]{subsection}

% parskip=half reserves 1em at the end of a paragraph; the full-width box needs
% that switched off locally.
\newcommand{\fillin}[2]{{\setlength{\parfillskip}{0pt plus 1fil}\noindent
  \fbox{\parbox{\dimexpr\linewidth-2\fboxsep-2\fboxrule\relax}{%
  \footnotesize\raggedright\textbf{#1} #2}}\par}}

\begin{document}

% ---------------------------------------------------------------- Deutsch
\fillin{Vor dem Versand ausfüllen:}{Name; Funktion; Telefon. Diesen
Kasten danach löschen.}

\begin{center}
  {\LARGE\bfseries Anspruch}\\[0.5ex]
  {\large Welche Pflegeleistungen einem Haushalt zustehen und was davon liegen bleibt}
\end{center}

Anspruch ist eine Webseite für Familien, die einen Angehörigen pflegen. Die Familie
beantwortet Fragen zum Alltag, und die Seite schätzt den Pflegegrad, den die
Begutachtung wahrscheinlich ergibt. Sie vergleicht die Leistungen, die daraus
folgen, mit dem, was der Haushalt schon bekommt, und nennt die nächsten Schritte mit
der gesetzlichen Grundlage.

\subsection*{Auf einen Blick}
\begin{itemize}
  \item \textbf{Sechs Sprachen} der Oberfläche: Deutsch, Englisch, Türkisch,
    Russisch, Polnisch, Arabisch. Die Fragen stehen auf Deutsch oder Englisch.
  \item \textbf{In der Regel 12 bis 20 Fragen} statt der 64 Fragen des
    amtlichen Instruments.
  \item \textbf{Ein Bereich, kein Mittelwert}, solange das Ergebnis nicht feststeht.
  \item \textbf{Gesetzliche Grundlage bei jedem Betrag.} Ein Betrag ohne Quelle
    wird nicht angezeigt.
  \item \textbf{Auf dem Gerät, ohne Konto.} Getippte Antworten bleiben im Browser
    der Familie.
  \item \textbf{Ein PDF für die Pflegekasse}, auf Deutsch oder Englisch, mit der
    Schätzung, den Beträgen, den Schritten und den gegebenen Antworten.
\end{itemize}

\subsection*{Was Sie als Organisation bekommen}
\begin{itemize}
  \item \textbf{Ihren eigenen Link.} Besuche darüber werden Ihrer Organisation
    zugeordnet, nicht einer Person.
  \item \textbf{Die Zahlen zu Ihrem Link}, sobald die Zählung eingeschaltet ist,
    per Link zu einer Übersicht oder als monatliche Zusammenstellung, nur als
    Summen.
  \item \textbf{Nichts Persönliches.} Keine Namen, keine Antworten, kein
    Pflegegrad. Familien brauchen kein Konto.
  \item \textbf{Einen Verweis auf die Beratung.} Die Seite sagt jeder Familie,
    dass ihr die kostenlose Pflegeberatung nach §~7a SGB~XI zusteht. Sie schickt
    Familien also zur Beratung, nicht von ihr weg.
\end{itemize}

\subsection*{Was Anspruch nicht ist}
\begin{itemize}
  \item Keine Begutachtung durch den Medizinischen Dienst. Es ist eine Schätzung
    aus den Angaben der Familie.
  \item Keine Rechtsberatung.
  \item Noch nicht von Muttersprachlern geprüft: die türkische, russische,
    polnische und arabische Oberfläche. Bisher hat nur ein Sprachmodell sie gegen
    den deutschen Text geprüft.
\end{itemize}
{\small Beträge und Punkteschwellen zuletzt geprüft am 22.10.2026. Das
Pflegeneuordnungsgesetz ist ein Kabinettsentwurf vom 30.09.2026; wird es wie
entworfen beschlossen, ändern sich die meisten Beträge ab dem 01.01.2027.\par}

\subsection*{Datenschutz}
Was Familien tippen, verlässt ihr Gerät nicht. Bei Antworten per Sprache erkennt der
Browser die Sprache, nicht die Seite. Wenn gezählt wird, dann ohne Cookies, und Sie
sehen nur Summen.

\subsection*{Kontakt}
\blank{Name}, \blank{Funktion} \quad kinlotangiri911@gmail.com \quad \blank{Telefon}

\newpage
% ---------------------------------------------------------------- English
\begin{otherlanguage}{english}
\fillin{Fill in before sending / Vor dem Versand ausfüllen:}{Name; Funktion
(role); Telefon (phone). The same blanks as on the German page. Delete
this box afterwards.}

\begin{center}
  {\LARGE\bfseries Anspruch}\\[0.5ex]
  {\large Which care benefits a household is entitled to, and what goes unclaimed}
\end{center}

Anspruch is a website for families who care for a relative. The family answers
questions about everyday life, and the site estimates the Pflegegrad the assessment
would likely arrive at. It compares the benefits that follow with what the household
already receives, and lists the next steps with their statutory basis.

\subsection*{At a glance}
\begin{itemize}
  \item \textbf{Six interface languages}: German, English, Turkish, Russian,
    Polish, Arabic. The questions are in German or English.
  \item \textbf{Typically 12 to 20 questions} instead of the official
    instrument's 64 questions.
  \item \textbf{A range, never a midpoint}, while the result is not settled.
  \item \textbf{Statutory basis on every figure.} A figure without a source is not
    shown.
  \item \textbf{On the device, no account.} Typed answers stay in the family's
    browser.
  \item \textbf{A PDF for the Pflegekasse}, in German or English, with the
    estimate, the amounts, the steps and the answers given.
\end{itemize}

\subsection*{What you get as an organisation}
\begin{itemize}
  \item \textbf{Your own link.} Visits through it are attributed to your
    organisation, not to a person.
  \item \textbf{The numbers for your link}, once counting is switched on, as a
    link to an overview or as a monthly summary, as totals only.
  \item \textbf{Nothing personal.} No names, no answers, no Pflegegrad. Families
    need no account.
  \item \textbf{A pointer to advice.} The site tells every family that free
    statutory care advice under § 7a SGB XI is their right. It sends families
    towards advice, not away from it.
\end{itemize}

\subsection*{What Anspruch is not}
\begin{itemize}
  \item Not an assessment by the Medizinischer Dienst. It is an estimate from what
    the family enters.
  \item Not legal advice.
  \item Not yet reviewed by native speakers: the Turkish, Russian, Polish and
    Arabic interface. So far only a language model has checked it against the
    German text.
\end{itemize}
{\small Amounts and points thresholds last checked on 22 Oct 2026. The
Pflegeneuordnungsgesetz is a cabinet bill of 30 Sep 2026; if it passes as drafted,
most amounts change from 1 Jan 2027.\par}

\subsection*{Privacy}
What families type does not leave their device. When answering by voice, the
browser does the speech recognition, not the site. Where counting is on, it uses no
cookies, and you see totals only.

\subsection*{Contact}
\blank{Name}, \blank{Funktion} \quad kinlotangiri911@gmail.com \quad \blank{Telefon}
\end{otherlanguage}

\end{document}
